% DOTE2011G · In-Class Exercise 1
% Mid-Autumn Festival lantern riddle (猜燈謎)
% Full house counting with combinations
%
% Recommended: xelatex Ex1.tex
% (ctexart needs XeLaTeX or LuaLaTeX for Chinese)

\documentclass[11pt,a4paper,UTF8]{ctexart}

\usepackage[margin=2.2cm]{geometry}
\usepackage{amsmath,amssymb}
\usepackage{enumitem}
\usepackage{setspace}
\usepackage{xcolor}
\usepackage{tcolorbox}
\tcbuselibrary{skins}
\usepackage{fancyhdr}
\usepackage{microtype}

\definecolor{night}{HTML}{0B1C3A}
\definecolor{lantern}{HTML}{C43C00}
\definecolor{paper}{HTML}{FFF8E8}

\pagestyle{fancy}
\fancyhf{}
\fancyhead[L]{\small DOTE2011G · Statistical Analysis for Business Decisions}
\fancyhead[R]{\small In-Class Exercise 1}
\fancyfoot[C]{\thepage}
\renewcommand{\headrulewidth}{0.4pt}

\setlist[enumerate]{itemsep=0.7em, topsep=0.4em}
\setlist[itemize]{itemsep=0.25em, topsep=0.2em}

\newcommand{\answerbox}[1][2.2cm]{%
  \par\vspace{0.35em}%
  \noindent
  \begin{tcolorbox}[
    colback=white,
    colframe=black!25,
    boxrule=0.6pt,
    arc=1mm,
    left=4pt,right=4pt,top=4pt,bottom=4pt,
    height=#1,
    valign=top
  ]
  \vspace{0.2em}
  \end{tcolorbox}%
}

\begin{document}

\begin{center}
  {\Large\bfseries\color{night} 中秋猜燈謎}\\[0.25em]
  {\large\bfseries Mid-Autumn Lantern Riddle}\\[0.55em]
  {\normalsize Counting a ``full house'' with combinations}\\[0.7em]
  {\small Participation exercise · please show your working}
\end{center}

\vspace{0.4em}

\begin{tcolorbox}[
  enhanced,
  colback=paper,
  colframe=lantern,
  boxrule=1.1pt,
  arc=2mm,
  left=10pt,right=10pt,top=8pt,bottom=8pt,
  title={\bfseries 謎面 · The riddle},
  fonttitle=\color{white}\bfseries,
  coltitle=white,
  colbacktitle=lantern
]
\begin{center}
  {\large
  三人同顏一對和，\\[0.15em]
  五張牌裡月婆娑。\\[0.15em]
  花燈若問幾許亮？\\[0.15em]
  請君細數莫猜錯。
  }
\end{center}

\vspace{0.55em}
\textit{Three of one face, and a matching pair --- five cards under the Mid-Autumn moon.
How many such lanterns can a full deck light?}
\end{tcolorbox}

\vspace{0.7em}

\noindent\textbf{Decode the riddle first.}
「三人同顏」means three cards of the \emph{same number}.
「一對和」means one \emph{pair} (two cards of another number).
Together this is a \textbf{full house}: three-of-a-kind $+$ one pair.

\vspace{0.55em}

\noindent\textbf{Deck rules (read carefully).}
A standard deck has \textbf{52 cards}: \textbf{4 suits} (kinds), and in each suit the numbers run from \textbf{1 to 13}.
So there are 13 ranks, and for each rank there are exactly 4 cards (one per suit).
You form a hand of \textbf{5 cards}. Order does \emph{not} matter --- count with combinations $C(n,k)$, not permutations.

\vspace{0.85em}
\noindent{\large\bfseries\color{night} 解謎任務 · Your tasks}

\begin{enumerate}[leftmargin=1.5em]
  \item \textbf{Pattern check.}
  In one sentence, state what hand the riddle is asking you to count
  (three of one rank $+$ a pair of another rank).
  \answerbox[1.6cm]

  \item \textbf{Build the count with combinations.}
  Write a product of combination symbols, then evaluate it, for the number of distinct full-house hands.

  Suggested steps (you may rearrange, but show the reasoning):
  \begin{itemize}
    \item choose the rank for the three-of-a-kind;
    \item choose which 3 of the 4 suits for that rank;
    \item choose the rank for the pair (different from the first);
    \item choose which 2 of the 4 suits for that pair.
  \end{itemize}
  \answerbox[4.2cm]

  \item \textbf{(Optional bonus)}
  There are $C(52,5)$ possible 5-card hands.
  Write the probability of a full house as a fraction
  (you need not fully simplify by hand if the numbers are large).
  \answerbox[2.0cm]
\end{enumerate}

\vspace{0.6em}
\noindent\rule{\textwidth}{0.4pt}\\[0.35em]
{\small
Name: \underline{\hspace{5.5cm}}
\quad SID / Group: \underline{\hspace{3.2cm}}
\quad Date: \underline{\hspace{2.8cm}}
}\\[0.55em]
{\small
Submit as instructed on Blackboard. This counts toward \textbf{course participation} (in-class exercises).
中秋快樂 · Happy Mid-Autumn Festival.
}

\end{document}
